% TOP_PROBLEM: {"id": 30006784, "title": "Nonvanishing conjecture for projective log canonical pairs", "category_id": 5, "category": {"id": 5, "name": "algebraic_geometry", "display_name": "Algebraic Geometry", "description": "Geometric objects defined by polynomial equations.", "slug": "algebraic-geometry", "order_index": 5, "created_at": "2026-07-31T15:26:25.671Z"}, "status": "open"}
% FIRST_POSTED: 2026-09-27
% !TeX program = lualatex
% Compile twice: lualatex open_problems_500.tex
% All references and prose are embedded; no BibTeX database or external inputs.
\documentclass[11pt,a4paper]{article}
\usepackage[margin=24mm,headheight=14pt]{geometry}
\usepackage{fontspec}
\setmainfont{Latin Modern Roman}
\setsansfont{Latin Modern Sans}
\setmonofont{Latin Modern Mono}
\usepackage{amsmath,amssymb,mathtools}
\usepackage{unicode-math}
\setmathfont{Latin Modern Math}
\usepackage{microtype}
\usepackage{enumitem,xurl,needspace}
\usepackage[dvipsnames]{xcolor}
\usepackage{hyperref}
\hypersetup{unicode=true,colorlinks=true,linkcolor=MidnightBlue,urlcolor=MidnightBlue,
 pdftitle={500 Mathematical Problem Statements},pdfauthor={Source-annotated compilation},bookmarksnumbered=true}
\usepackage{bookmark}
\usepackage{tocloft}
\setlength{\cftsecnumwidth}{3em}
\cftsetpnumwidth{2.5em}
\cftsetrmarg{4em}
\usepackage{titlesec}
\titleformat{\section}{\Large\bfseries\raggedright}{\thesection}{0.7em}{}
\usepackage{fancyhdr}
\pagestyle{fancy}\fancyhf{}
\fancyhead[L]{\small\sffamily Mathematical problem statements}
\fancyhead[R]{\small Source edition 22}
\fancyfoot[C]{\thepage}
\setlength{\parindent}{0pt}
\setlength{\parskip}{5pt plus 1pt}
\setlength{\emergencystretch}{3em}
\setcounter{tocdepth}{1}
\setcounter{secnumdepth}{1}
\setlist[itemize]{leftmargin=3em,itemsep=5pt,topsep=4pt,parsep=0pt}
\allowdisplaybreaks
\newcommand{\N}{\mathbb{N}}
\newcommand{\Z}{\mathbb{Z}}
\newcommand{\Q}{\mathbb{Q}}
\newcommand{\R}{\mathbb{R}}
\newcommand{\C}{\mathbb{C}}
\newcommand{\F}{\mathbb{F}}
\newcommand{\A}{\mathbb{A}}
\newcommand{\PP}{\mathbb P}
\newcommand{\E}{\mathbb{E}}
\newcommand{\eps}{\varepsilon}
\newcommand{\dd}{\,\mathrm d}
\newcommand{\sref}[2]{\hyperlink{source:#1:#2}{[S#2]}}
\newcommand{\eref}[2]{\hyperlink{extra:#1:#2}{[E#2]}}
% Turn the next switch off for a shorter reading copy. The quotations preserve
% every exactTarget field, including qualifications omitted from a short summary.
\newif\ifshowcatalogue
\showcataloguetrue
\newcommand{\cataloguescope}[1]{\ifshowcatalogue
 \paragraph{Exact scope in the supplied catalogue.}
 \begin{quote}\small #1\end{quote}\fi}

% Compile twice with: lualatex 30006784.tex
\numberwithin{equation}{section}
\hypersetup{pdftitle={Research attempt -- problem 30006784},pdfauthor={Research notebook; original compilation retained}}
\fancyhead[L]{\small\sffamily Research notebook}
\fancyhead[R]{\small\sffamily Problem 30006784 | 27 September 2026}
\begin{document}
\setcounter{section}{0}
% ================================================================
% problemId: 30006784
% Canonical problem ID: 30006784
% exactTarget SHA-256: 531f036f8babea440fd4e288732047e75f37c2ed4e7e9b11ab1b3db57c1bfded
\clearpage
\section*{\texorpdfstring{Nonvanishing conjecture for projective log canonical pairs}{Nonvanishing conjecture for projective log canonical pairs}}\label{30006784}
\subsection{Definitions and mathematical statement}
A log canonical pair $(X,\Delta)$ has normal $X$, effective rational boundary $\Delta$, $K_X+\Delta$ rational Cartier, and all discrepancies at least $-1$ on resolutions. A divisor is pseudo-effective if its numerical class lies in the closure of the cone generated by effective divisors. For projective complex pairs, the assertion is
\[
 K_X+\Delta\text{ pseudo-effective}
 \quad\Longrightarrow\quad
 H^0(X,\mathcal O_X(m(K_X+\Delta)))\ne0
\]
for some sufficiently divisible $m>0$ making that divisor integral Cartier. Equivalently its Iitaka dimension is nonnegative. Numerical approximation by effective classes is not itself the existence of a section or linear equivalence to an effective divisor; that is precisely the missing implication.
\par\smallskip\textit{Formulation and concept references:} \sref{30006784}{1}.
\cataloguescope{Let (X, Delta) be a projective log canonical pair over the complex numbers, with Delta a Q-divisor and K\_X+Delta Q-Cartier. If K\_X+Delta is pseudo-effective, prove that kappa(K\_X+Delta) >= 0; equivalently, some sufficiently divisible positive multiple of K\_X+Delta is linearly equivalent to an effective divisor.}
\subsection{Short English statement}
Must a pseudo-effective log canonical divisor on a projective complex pair acquire a nonzero section after some positive integral multiple?
% BEGIN RESEARCH ADDITION ATTEMPT-194
\subsection{Research attempt}

\paragraph{Current frontier.}
The original source studies direct images and special cases of nonvanishing; it is not used here as a proof for every log canonical pair \sref{30006784}{1}. Lazi\'c--Meng prove a dimension-inductive nonvanishing result for uniruled log canonical pairs and obtain the four-dimensional uniruled case from lower-dimensional results \eref{30006784}{1}, Theorem~1.1 and Corollary~1.2. The general obstacle in the printed target is the passage from a numerical limiting condition to actual sections. We examine the residual obstruction in $\operatorname{Pic}^0$, rather than assuming numerical and linear equivalence coincide.

\paragraph{Attack route.}
One route perturbs $K_X+\Delta$ by an ample divisor and tries to pass sections to a limit; it would require a uniform control on their multiples and twisting classes. A second uses adjunction to a boundary component, but needs surjectivity of restriction of pluricanonical sections. A third first produces numerical effectivity and then removes a topologically trivial twist. We pursue the third because it admits an exact torsion-point criterion and an explicit stress example. The criterion will not be asserted for arbitrary line bundles merely because it is desirable for adjoint ones.

\paragraph{Attempt.}
Put $D=K_X+\Delta$. For an integer $m>0$ making $mD$ Cartier, set $A=\mathcal O_X(mD)$ and define
\[
 V^0(A)=\{\alpha\in\operatorname{Pic}^0(X):H^0(X,A\otimes\alpha)\ne0\}.
\]
This definition is valid for the projective normal variety in the target. We make no smoothness assumption on $X$ in the following elementary argument.

\textbf{Torsion-point extraction lemma.} If $V^0(A)$ contains a torsion point $\alpha$ of order $r$, then $H^0(X,A^{\otimes r})\ne0$.
Choose a nonzero section $s$ of $A\otimes\alpha$. Its tensor power $s^{\otimes r}$ is a nonzero section of $A^{\otimes r}\otimes\alpha^{\otimes r}\simeq A^{\otimes r}$. Nonzero follows by evaluating at the generic point of the integral variety. Thus $rmD$ is linearly equivalent to an effective divisor, exactly as required in the original conclusion.

A concrete sufficient pair of missing inputs is now the following, for each pseudo-effective adjoint divisor in the question:
\begin{enumerate}[label=(\roman*),leftmargin=2.5em]
\item For some admissible $m$, there are an effective Cartier divisor $E$ and $L\in\operatorname{Pic}^0(X)$ with $A\simeq\mathcal O_X(E)\otimes L$.
\item For this $A$, the nonempty locus $V^0(A)$ contains a torsion point; a sufficient stronger property is that one of its irreducible components is a torsion translate of an abelian subvariety.
\end{enumerate}
Input (i) puts $L^{-1}$ in $V^0(A)$. Input (ii) and the lemma finish the target. If one starts instead with a merely numerically trivial difference, the additional passage to an algebraically trivial multiple must also be justified; (i) is deliberately stated in the precise form actually used. We have not proved either input in the generality of the conjecture.

The condition in (ii) is substantially stronger than nonemptiness. To test it, let $C$ be a complex elliptic curve and $L\in\operatorname{Pic}^0(C)$ a non-torsion line bundle. For every $m\ge1$,
\begin{equation}\label{eq194:elliptic}
 H^0(C,L^{\otimes m})=0.
\end{equation}
Indeed, any nonzero section defines an effective divisor of degree zero, hence the zero divisor, and trivializes $L^{\otimes m}$, contradicting non-torsion. Nevertheless $L$ is numerically trivial, hence pseudo-effective as a numerical class, and
\[
 V^0(L)=\{L^{-1}\}.
\]
This is a nonempty locus with no torsion point. It disproves the proposed numerical-to-linear argument for unrestricted divisors. It does \emph{not} refute the nonvanishing conjecture: on an elliptic curve $K_C$ is trivial, and an effective boundary of degree zero must be zero, so this non-torsion line bundle cannot be the stated adjoint divisor.

One might try to remove $L$ by a finite cover. That does not work formally either. If $f:Y\to X$ is finite flat of degree $d$ and $f^*L$ is trivial, the norm identity gives
\[
 L^{\otimes d}\simeq\operatorname{Nm}_f(f^*L)\simeq\mathcal O_X.
\]
Thus finite covers cannot trivialize a genuinely non-torsion element of $\operatorname{Pic}^0$ by this mechanism. More positively, if a nonzero section of $f^*A$ is produced on an integral $Y$, its norm is a nonzero section of $A^{\otimes d}$. A finite-cover argument therefore can descend nonvanishing, but must supply an actual upstairs section rather than just numerical triviality.

The other attack routes have comparably explicit gaps. Sections of $m(D+\varepsilon H)$ with $m$ depending on $\varepsilon$ do not lie in one fixed finite-dimensional vector space, so compactness of a projective space of sections does not apply. For a Cartier boundary component $S$, the elementary restriction sequence
\[
 0\longrightarrow\mathcal O_X(mD-S)\longrightarrow\mathcal O_X(mD)
 \longrightarrow\mathcal O_S(mD|_S)\longrightarrow0
\]
shows that an effective restricted adjoint class lifts only if its connecting image in $H^1(X,\mathcal O_X(mD-S))$ vanishes. Pseudo-effectivity alone supplies no such vanishing statement. This explains why neither a limiting argument nor adjunction by itself fills input (ii).

\paragraph{Stress test.}
The torsion lemma works for a singular integral projective $X$ and requires no analytic continuation of sections. It retains rational coefficients by passing to admissible integral multiples. It does not mistake numerical approximation for an effective representative, and no general torsion-translate theorem for arbitrary $V^0(A)$ is invoked: \eqref{eq194:elliptic} would contradict one. Any application of a support-locus theorem for pluricanonical bundles must check its precise singularity, boundary, and morphism hypotheses against a general log canonical pair. The cited special-case results do not automatically supply that check in every dimension.

\paragraph{Outcome.}
\textbf{Outcome: Conditional reduction.}
The target follows from effective representation up to a $\operatorname{Pic}^0$ twist together with a torsion-point property for the resulting support locus. The extraction and norm steps are proved, and the non-torsion elliptic example demonstrates exactly why purely numerical reasoning is insufficient. The adjoint-specific input that would force a torsion twist, and the required numerical effectivity in general, remain unproved.

% END RESEARCH ADDITION ATTEMPT-194
\subsection{Sources}
\begin{itemize}\raggedright
\item[\textnormal{[S1]}]\hypertarget{source:30006784:1}{} Houari Benammar Ammar, Remarks on the Direct Image Sheaf of Logarithmic Pluricanonical Bundles and the Non-Vanishing Conjecture, arXiv:2501.04232v1 (2025).
\url{https://arxiv.org/html/2501.04232v1}.
{\small\textit{Catalogue formulation source.}}
% BEGIN RESEARCH ADDITION REFERENCES-194
\item[\textnormal{[E1]}]\hypertarget{extra:30006784:1}{} Vladimir Lazi\'c and Fanjun Meng, \emph{On Nonvanishing for uniruled log canonical pairs}, arXiv:1907.11991v2; Electronic Research Archive 29 (2021). Theorem 1.1 and Corollary 1.2 distinguish the inductive uniruled results from full nonvanishing.
\url{https://arxiv.org/html/1907.11991v2}.
{\small\textit{Research reference; consulted 27 September 2026.}}
% END RESEARCH ADDITION REFERENCES-194
\end{itemize}

\end{document}
